\section*{附录：关键引语索引}

访谈里最值得反复回看的内容，往往并不是某个新闻点，而是这些被谢赛宁反复压缩成一句话的判断。它们之所以有力量，不是因为措辞漂亮，而是因为每一句背后都对应着一个长期问题意识。下面把全场最重要的一组引语与其语境整理成索引，方便后续回看时快速定位其思想位置。

\begin{longtable}{>{\raggedright\arraybackslash}p{0.34\textwidth} >{\raggedright\arraybackslash}p{0.58\textwidth}}
\toprule
关键引语 & 所对应的问题意识 \\
\midrule
\endfirsthead
\toprule
关键引语 & 所对应的问题意识 \\
\midrule
\endhead
\textit{``眼睛是唯一一个暴露在真实世界里面的大脑部分...解决视觉不是要解决视觉本身，而是要解决智能本身''} & 这是整场访谈最核心的理论前提。它把视觉从具体任务提升为智能与现实世界之间的接口，也解释了为什么谢赛宁后来会一路走向 world model。 \\
\textit{``Research 从来不是一个线性的发展，一个线性发展的 research 永远不是好的 research''} & 这句话定义了他对创造性工作的基本理解。好的研究必须包含偏航、冲突和重写问题的过程，否则说明问题本身还不够深。 \\
\textit{``一开始你想的 idea 不是你的 idea，探索中的 idea 才是属于你的 idea''} & 它强调真正属于研究者的洞见只能在做事过程中长出来，而不是在空想阶段一次性想明白。探索不是准备动作，而是 idea 生成本身。 \\
\textit{``凡所有像皆是虚妄，若见诸像非相，即见如来''} & 在访谈语境里，这不是宗教引文展示，而是对 research taste 的描述：不要被论文表象、术语包装和局部指标迷惑，要追问背后的实质变量。 \\
\textit{``LLM 终将凋零...老兵不死，终将凋零''} & 这句判断的重点不是否定语言模型，而是否定把 LLM 当作终局。它指向的，是语言模型在真实世界表征与规划能力上的结构性局限。 \\
\textit{``过去是 download internet 的时代，现在是 download human 的时代''} & 它重新定义了下一阶段 AI 数据基础。文本互联网不再是唯一中心，视频、感知、行为和人类在世界中的轨迹将成为更关键的训练来源。 \\
\textit{``能够重新造出一只松鼠，要比人类文明在 530 million years 最后 8 秒创造的东西伟大得多''} & 这句话用动物智能批评狭义 AGI 叙事，提醒我们不要把高分符号任务误当成智能本身。真实世界中的生存、感知与行动，难度远高于若干抽象 benchmark。 \\
\textit{``I am not the special one, I am the normal one''} & 这既是对个人神话的拒绝，也是对“电池精神”的肯定。研究共同体需要的不是少数被神化的天才，而是愿意长期给问题和他人供能的人。 \\
\textit{``不要在乎一个 point estimate...所有的评价最后都会是一个积分''} & 它总结了谢赛宁对评审、中稿、奖项与短期成败的态度。局部评价是高噪声抽样，真正重要的是更长时间跨度上的累计影响。 \\
\textit{``每个人都是这个世界上一个变量，有可能你就是最重要的那个变量''} & 这是他对年轻研究者最直接的鼓励之一。面对高度竞争和强路径依赖的系统，个人仍然应该保留主动进入场域、主动联系机会、主动承担风险的意志。 \\
\bottomrule
\end{longtable}

如果把这十句引语并排来看，会发现它们其实都指向同一个中心命题：真正值得投入一生的问题，通常不在短期评价最清晰的地方，而在那些仍然需要你亲自进入、亲自摸索、亲自承担不确定性的地方。谢赛宁无论谈视觉、研究方法、创业还是人生，都没有离开这个命题。

\end{document}
